\section{引言：大模型时代的强化学习浪潮}

2025年，强化学习（RL）在大语言模型（LLM）领域掀起了一场技术革命。从年初DeepSeek R1的发布开始，各家旗舰模型的训练都离不开强化学习。在2025年"青稞"AI嘉年华的RL专题讨论中，四位在强化学习领域经验丰富的嘉宾——崔干曲（上海人工智能实验室）、胡建（OpenRLHF项目发起人，Emilia Research）、郑楚杰（通义千问团队）和李英儒（NovaStar Research）——围绕强化学习算法设计、模型架构适配、训练框架演进、训练稳定性以及未来发展等核心话题，进行了深入的分享与讨论。

\begin{importantbox}{本次讨论的核心议题}
\begin{enumerate}
    \item 强化学习算法设计的出发点与心路历程
    \item 模型架构如何适配强化学习训练
    \item RL训练框架的现状与未来
    \item Entropy、Mismatch等关键指标与训练稳定性
    \item 2026年强化学习的发展方向展望
\end{enumerate}
\end{importantbox}

\subsection{嘉宾介绍与2025年关键进展}

\textbf{崔干曲}（上海人工智能实验室，青年科学家）：毕业于清华大学，主要研究强化学习在大模型推理能力提升上的算法与机理。代表工作包括PRIME（复刻O1的开创性工作）、Implicit PRM（隐式过程奖励模型）以及首个达到IPHO金牌水平的开源物理推理模型P1。

\textbf{胡建}（网名"初期"，OpenRLHF项目发起人）：现任职于Emilia Research，长期处于RL的infra与算法联合优化的交叉领域。提出了Reinforce++等有影响力的算法，并主导了OpenRLHF框架的设计与开发。

\textbf{郑楚杰}（通义千问团队）：专注于大模型RL训练的稳定性研究，核心关注如何让RL训练长期稳定地持续下去。提出了GSPO算法，并应用于千问2507等模型的发版训练。

\textbf{李英儒}（NovaStar Research）：强化学习理论算法出身，关注训退不一致（Train-Inference Mismatch）问题以及RL优化理论在大模型下的适配。致力于从第一性原理出发，为现有的RL recipe提供理论解释与指导。

\begin{knowledgebox}{2025年RL领域的共识与趋势}
四位嘉宾不约而同地指出：2025年最有影响力的进展大多不在算法层面，而在\textbf{Infra}（基础设施）和\textbf{数据}层面。RL算法本身仍然以PPO（2017年）甚至Policy Gradient（提出数十年）为基础，但如何在现代GPU上实现高效稳定的大规模训练，这方面的创新和突破是巨大的。代表性进展包括：
\begin{itemize}
    \item \textbf{训退不一致（Train-Inference Mismatch）}：被广泛认识为影响训练稳定性的关键问题
    \item \textbf{MoE模型的RL稳定性}：Mixture-of-Experts架构给RL训练带来了特有的挑战
    \item \textbf{Hybrid架构}：超长上下文与Agent场景下的RL稳定性
    \item \textbf{Importance Sampling}：Truncated IS等技术从infra观察中产生算法改进
\end{itemize}
\end{knowledgebox}

\subsection{本章小结}

2025年是大模型RL的"文艺复兴"年——理论基础早已存在，但在大模型的特殊场景下进行适配和工程化落地才是真正的创新前沿。Infra与算法的联合优化成为主流范式，纯粹的理论推导已不足以驱动进展，必须结合底层系统的观察来指导算法改进。


